% Approved shared presentation for the full library and abstract preview.
% 全库阅读版与抽象范畴预览共用的已批准展示样式。
\setmainfont{texgyrepagella-regular.otf}[
  BoldFont=texgyrepagella-bold.otf,
  ItalicFont=texgyrepagella-italic.otf,
  BoldItalicFont=texgyrepagella-bolditalic.otf]
\setsansfont{texgyreheros-regular.otf}[BoldFont=texgyreheros-bold.otf]
\usepackage[most]{tcolorbox}
\usepackage{titlesec}
\usepackage{fancyhdr}
\definecolor{catsBlue}{HTML}{263B53}
\definecolor{catsInk}{HTML}{253342}
\definecolor{catsMuted}{HTML}{647487}
\definecolor{catsPaper}{HTML}{F7F9FC}
\definecolor{catsBorder}{HTML}{DCE4ED}
\geometry{left=22mm,right=22mm,top=23mm,bottom=23mm,headheight=16pt}
\linespread{1.08}
\setlength{\parskip}{.6em plus .1em minus .1em}
% Allow intact inline identifiers to move to the next line in justified prose.
% 两端对齐的正文允许完整的行内标识符移至下一行。
\setlength{\emergencystretch}{10em}
\titleformat{\chapter}[display]
  {\color{catsBlue}\normalfont\LARGE\bfseries}
  {\sffamily\small ROCQ-CATS}{6pt}{}
\titlespacing*{\chapter}{0pt}{0pt}{16pt}
\titleformat{\section}{\color{catsBlue}\Large\bfseries}{}{0pt}{}
\titleformat{\subsection}{\color{catsBlue}\large\bfseries}{}{0pt}{}
\pagestyle{fancy}
\fancyhf{}
\fancyhead[L]{\sffamily\footnotesize\color{catsMuted}Rocq-Cats}
\fancyhead[R]{\sffamily\footnotesize\color{catsMuted}\nouppercase{\leftmark}}
\fancyfoot[R]{\sffamily\small\color{catsBlue}\thepage}
\renewcommand{\headrulewidth}{.3pt}
\renewcommand{\chaptermark}[1]{\markboth{#1}{}}
\fancypagestyle{plain}{\fancyhf{}
  \fancyfoot[R]{\sffamily\small\color{catsBlue}\thepage}
  \renewcommand{\headrulewidth}{0pt}}
% Unstarred break padding keeps continued frames within the page height.
% 不带星号的跨页留白保持代码框位于页面高度以内，不自动扩大页面。
\tcbset{cats code/.style={enhanced jigsaw,breakable,
  colback=catsPaper,colframe=catsBorder,coltext=catsInk,
  boxrule=.5pt,arc=5pt,outer arc=5pt,
  left=9pt,right=9pt,top=7pt,bottom=7pt,
  before skip=9pt,after skip=9pt,
  colbacktitle=catsPaper,coltitle=catsBlue,
  fonttitle=\sffamily\footnotesize\bfseries,
  titlerule=.4pt,toptitle=5pt,bottomtitle=5pt,
  pad at break=5pt,lines before break=3}}
\tcbset{cats semantic/.style={},
  cats theorem/.style={colback=white,
    borderline west={1pt}{0pt}{catsBlue}}}
\newcommand{\catsCodeTitle}{Rocq}
\newcommand{\catsCodeKind}{code}
\newcommand{\catsContinued}{continued}
\newenvironment{catscode}[2]{%
  \begin{tcolorbox}[cats code,cats semantic,
    title={#1\hfill\normalfont\scriptsize\ttfamily #2},
    title after break={#1\ (\catsContinued)\hfill
      \normalfont\scriptsize\ttfamily #2}]
  \small\ttfamily\raggedright\setlength{\parskip}{0pt}
  \setlength{\baselineskip}{12pt}\color{catsInk}}
  {\end{tcolorbox}}
\renewenvironment{coqdoccode}
  {\begin{catscode}{\catsCodeTitle}{\catsCodeKind}}
  {\end{catscode}}
\renewcommand{\coqdockw}[1]{\textcolor{catsBlue}{\textbf{#1}}}
\renewcommand{\coqdocvar}[1]{\texttt{#1}}
\renewcommand{\coqdoccst}[1]{\texttt{#1}}
\renewcommand{\coqdocmod}[1]{\texttt{#1}}
\renewcommand{\coqdocax}[1]{\texttt{#1}}
\renewcommand{\coqdocind}[1]{\texttt{#1}}
\renewcommand{\coqdocconstr}[1]{\texttt{#1}}
\renewenvironment{coqdoccomment}{\begingroup\color{catsMuted}\ttfamily(*}
  {*)\endgroup}
\renewcommand{\coqdocemptyline}{\par\vspace{3pt}}
